\documentclass[11pt]{article}
\usepackage[margin=1in]{geometry}
\usepackage[T1]{fontenc}
\usepackage[utf8]{inputenc}
\usepackage{lmodern}
\usepackage{microtype}
\usepackage{amsmath,amssymb}
\usepackage{booktabs}
\usepackage{graphicx}
\usepackage{hyperref}
\usepackage{enumitem}
\usepackage{tabularx}

\title{Math4AI Capstone Report Title}
\author{Team Member 1 \and Team Member 2 \and Team Member 3}
\date{}

\begin{document}
\maketitle

\begin{abstract}
State the central question, the models compared, the main experimental evidence, and the conclusion in 120--180 words.
\end{abstract}

\section{Introduction}
Frame the capstone question and explain why comparing a linear model to a one-hidden-layer neural network is scientifically meaningful.

\section{Background and Mathematical Setup}
This section summarizes the mathematical setup used throughout the report.

\subsection{Backpropagation Concept (No Equations)}
Backpropagation is an efficient bookkeeping method for the chain rule. In the forward pass, the model computes intermediate quantities in sequence: hidden pre-activations, hidden activations, output scores, probabilities, and loss. In the backward pass, sensitivity to the loss is propagated in reverse through the same computation graph. Each local derivative is simple; the full parameter gradients are obtained by repeatedly applying the chain rule.

The key practical point is efficiency. Reverse-mode differentiation computes gradients for all parameters with cost on the same order as one forward pass (up to a small constant factor), which is why deep models are trainable in practice. Conceptually, backprop answers: if the loss changes a little, which parameters are responsible, in what direction, and by how much?

\subsection{NLL Derivation for Softmax Cross-Entropy}
For one sample $(x_i,y_i)$ and class scores $s_i \in \mathbb{R}^k$, define
\[
p_{ij} = \frac{e^{s_{ij}}}{\sum_{\ell=1}^k e^{s_{i\ell}}}
      = \mathbb{P}(Y=j \mid X=x_i; \theta).
\]
If labels are categorical, the likelihood for one sample is
\[
\mathcal{L}_i(\theta) = p_{i,y_i}.
\]
Hence the dataset log-likelihood is
\[
\log \mathcal{L}(\theta)=\sum_{i=1}^n \log p_{i,y_i}.
\]
Maximizing average log-likelihood is equivalent to minimizing
\[
-\frac{1}{n}\sum_{i=1}^n \log p_{i,y_i},
\]
which is exactly the mean cross-entropy loss. Therefore, minimizing softmax cross-entropy is the same optimization objective as maximum likelihood estimation for a conditional categorical model.

\subsection{Complete Gradient Derivation (One-Hidden-Layer NN)}
Let batch input be $X\in\mathbb{R}^{n\times d}$ and one-hot labels $Y\in\mathbb{R}^{n\times k}$. With one hidden layer:
\[
Z_1 = XW_1^\top + \mathbf{1}b_1^\top,\quad
H = \tanh(Z_1),\quad
S = HW_2^\top + \mathbf{1}b_2^\top,\quad
P=\mathrm{softmax}_{\text{row}}(S).
\]
Use mean cross-entropy plus L2 on weights only:
\[
J = -\frac{1}{n}\sum_{i=1}^n\sum_{j=1}^k Y_{ij}\log P_{ij}
    + \frac{\lambda}{2}\left(\|W_1\|_F^2+\|W_2\|_F^2\right).
\]
For the output layer:
\[
\frac{\partial J}{\partial S}=\frac{1}{n}(P-Y),
\quad
\frac{\partial J}{\partial W_2}
=\left(\frac{\partial J}{\partial S}\right)^\top H + \lambda W_2,
\quad
\frac{\partial J}{\partial b_2}
=\left(\frac{\partial J}{\partial S}\right)^\top\mathbf{1}.
\]
Propagate to hidden pre-activations:
\[
\frac{\partial J}{\partial Z_1}
=\left(\frac{\partial J}{\partial S}\right)W_2
\odot \left(1-H\odot H\right),
\]
where $1-H\odot H$ is the elementwise derivative of $\tanh$.
Then
\[
\frac{\partial J}{\partial W_1}
=\left(\frac{\partial J}{\partial Z_1}\right)^\top X + \lambda W_1,
\quad
\frac{\partial J}{\partial b_1}
=\left(\frac{\partial J}{\partial Z_1}\right)^\top\mathbf{1}.
\]
L2 is applied to weights only; there is no regularization term on $b_1,b_2$.

\subsection{Geometric Interpretation (Gaussian vs Moons)}
In the linear Gaussian task, classes are generated from roughly affine-separable structure in input space. A softmax linear classifier creates hyperplane decision boundaries, which is well matched to this geometry. Therefore, a linear baseline is expected to perform strongly and often does not require additional representation depth to capture the dominant class structure.

In the moons task, classes lie on interleaving curved manifolds. A single affine boundary cannot separate the classes cleanly. The hidden layer introduces a nonlinear feature map (through $\tanh$), which can bend and re-embed the input geometry before the final linear readout. This model pairing is therefore meaningful: softmax tests whether linear structure is sufficient, while the one-hidden-layer network tests whether nonlinear representation is needed for the observed geometry.

\section{Methods}
Describe the datasets, preprocessing, softmax regression, neural network, training protocol, optimizer choices, and evaluation metrics. Make the digits protocol explicit and state exactly how validation and test data were used.

\section{Implementation Sanity Checks}
Briefly document at least three verification checks, such as a gradient sanity check, probability-sum check, tiny-subset loss decrease, tiny-subset overfitting test, or NaN/Inf check.

\section{Experiments}
Report the required core comparisons and ablations. Prefer a small number of well-designed figures over many shallow plots. Treat the following as questions your evidence must answer: when is the linear model sufficient, what representational change the hidden layer provides, when extra complexity fails to justify itself, what the failure case teaches, and what repeated-seed statistics add beyond a single run.

\section{Advanced Track}
Include exactly one of the following:
\begin{itemize}[leftmargin=1.4em]
    \item Track A: PCA/SVD and input geometry
    \item Track B: Prediction confidence and reliability
\end{itemize}

\section{Discussion}
Interpret the results. State clearly where linear structure was sufficient, where nonlinearity helped, and why. Distinguish what is supported by evidence from what remains uncertain. Use consistent notation and define symbols before using them.

\section{Limitations}
Explain what the project does not test and which conclusions should not be overgeneralized.

\section*{References}
Use a consistent citation style.

\end{document}
